% ============================================================
%  hcm_paper_manuscript.tex   --   FINAL COMPLETE MANUSCRIPT
%  Deciphering the Molecular Grammar of Hypertrophic
%  Cardiomyopathy: A Zero-Leakage Language Modeling
%  Approach to Sarcomeric Variant Pathogenicity
%
%  Format  : IEEEtran journal,compsoc (double-column)
%  Authors : RVCE Experiential Learning Project 2026-26
%
%  Part I  (Sec I–III-B) : teammate contribution, copied verbatim
%  Part II (Sec III-C+)  : second-author contribution
% ============================================================

\documentclass[conference]{IEEEtran}

%  -- - Packages  --  --  --  --  --  --  --  --  --  --  --  --  --  --  --  -- 
\usepackage[T1]{fontenc}
\usepackage{cite}
\usepackage{amsmath,amssymb,amsfonts}
\usepackage{algorithmic}
\usepackage{graphicx}
\usepackage{url}
\usepackage{booktabs}
\usepackage{multirow}
\usepackage{xcolor}
\usepackage{hyperref}
% IEEE standard prefers no colored hyperlinks for print compatibility
\hypersetup{
  hidelinks
}
%  -- - Custom commands  --  --  --  --  --  --  --  --  --  --  --  --  -- --
\newcommand{\Deltaemb}{\Delta\mathbf{e}}
\newcommand{\xtab}{\mathbf{x}^{\mathrm{tab}}}
\newcommand{\xesm}{\mathbf{x}^{\mathrm{esm}}}
\newcommand{\ztab}{\mathbf{z}^{\mathrm{tab}}}
\newcommand{\zesm}{\mathbf{z}^{\mathrm{esm}}}
\newcommand{\zfus}{\mathbf{z}^{\mathrm{fus}}}
\newcommand{\phat}{\hat{p}}
% ==============================================================
\begin{document}

%  -- - Title & Authors  --  --  --  --  --  --  --  --  --  --  --  --  -- --
\title{Deciphering the Molecular Grammar of Hypertrophic
Cardiomyopathy: A Zero-Leakage Language Modeling Approach
to Sarcomeric Variant Pathogenicity}

\author{%
  \IEEEauthorblockN{Computational Genomics Team}\\
  \IEEEauthorblockA{Department of Computer Science and Engineering,\\
  RV College of Engineering, Bengaluru, India 560059}
}

\maketitle

% ==============================================================
\begin{abstract}
Hypertrophic cardiomyopathy (HCM) is the most prevalent inherited
disorder of cardiac muscle, responsible for a disproportionate
fraction of sudden cardiac deaths in young individuals. Genomic
sequencing routinely reveals missense substitutions in nine
sarcomeric genes for which clinical pathogenicity classification
remains unresolved -- a bottleneck that directly impedes cascade
screening and risk stratification. We present a dual-model,
zero-leakage pathogenicity prediction framework that eliminates two
distinct leakage channels simultaneously: feature-level leakage is
prevented by excluding all ACMG/AMP-derivable covariates from the
training matrix, and homology leakage is prevented by evaluating
under Leave-One-Gene-Out (LOGO) cross-validation, which forces
out-of-gene inference under distribution shift. The framework trains
on 1,954 high-confidence labeled missense variants (1,441
pathogenic, 513 benign) across nine sarcomeric genes and deploys
on 4,524 unresolved Variants of Uncertain Significance. Two
complementary classifiers are developed: a Baseline Random Forest
operating on 51 hand-crafted biophysical descriptors, and a
Two-Tower Hybrid Neural Network that additionally fuses
1,280-dimensional differential evolutionary embeddings extracted
from a frozen 650M-parameter ESM-2 protein language model. Under
LOGO evaluation, the Baseline RF achieves a mean AUPRC of
\textbf{0.886} (95\% CI across all genes), outperforming all five
external benchmarks: EVE~(0.791), AlphaMissense~(0.800),
REVEL~(0.842), MetaRNN~(0.849), and CardioBoost~(0.876), across
seven of nine target genes. Post-hoc Platt scaling calibrates
raw neural probabilities to clinically meaningful likelihoods
(Brier score 0.145 on held-out \textit{TNNT2}). Residue-level
explainability is provided through In Silico Mutagenesis (ISM)
heatmaps that map predicted vulnerability to known biomechanical
interfaces -- including the myosin converter domain and actin-binding
cleft -- without explicit 3-D coordinate supervision. These results
demonstrate that disease-specific, zero-leakage representation
learning captures the sarcomeric molecular grammar and constitutes
a validated triage mechanism for clinical VUS re-stratification.
\end{abstract}

\begin{IEEEkeywords}
Hypertrophic cardiomyopathy, variant pathogenicity, protein language
model, ESM-2, two-tower neural network, random forest, missense
variants, data leakage, Leave-One-Gene-Out, Platt scaling, in silico
mutagenesis, VUS re-stratification, ClinVar, gnomAD.
\end{IEEEkeywords}


% ==============================================================
% SECTION I  --  INTRODUCTION
% (Teammate contribution  --  verbatim from Part1.tex)
% ==============================================================
\section{Introduction}

\subsection{Proteins, Mutations, and Variant Interpretation}

Proteins are the primary functional molecules of the cell, encoded
by the genome and assembled from sequences of amino acids that fold
into precise three-dimensional structures. This architecture
determines a protein's stability, binding partners, enzymatic
activity, and structural role within macromolecular machines. A
single missense mutation -- the substitution of one amino acid for
another -- can perturb this structure in ways ranging from
imperceptible to catastrophic, depending on the physicochemical
nature of the change and its location within the protein's
functional architecture. With modern next-generation sequencing now
capable of surveying the full coding content of the human genome,
clinical and research pipelines routinely surface thousands of such
variants per individual, the vast majority of whose functional
consequences remain uncharacterised.

Experimental approaches such as enzymatic assays and high-throughput
deep mutational scanning \cite{fowler2014deep} provide direct
functional evidence but cannot scale to the volume of variants
uncovered by population-level sequencing. Computational
\textit{in silico} tools have therefore become an essential complement,
leveraging evolutionary conservation, protein structure, and
statistical learning to assign pathogenicity likelihoods at scale.
These tools operate within the five-tier classification framework
defined by the American College of Medical Genetics and Genomics and
the Association for Molecular Pathology (ACMG/AMP)
\cite{richards2015} -- Pathogenic, Likely Pathogenic, Variant of
Uncertain Significance (VUS), Likely Benign, and Benign -- which
forms the operational target for any predictive model intended for
clinical deployment. Despite steady methodological advances, no
existing tool has resolved the challenge of disease-specific variant
interpretation in clinically actionable gene panels, a gap that
motivates the present work.

\subsection{Hypertrophic Cardiomyopathy: Clinical and Genetic Context}

Hypertrophic cardiomyopathy (HCM) is the most prevalent inherited
disorder of the cardiac muscle, affecting approximately one in 200
to 500 individuals in the general population and following an
autosomal dominant inheritance pattern \cite{maron2013}. The
condition is defined by unexplained left ventricular hypertrophy
(LVH) in the absence of another systemic or cardiac disease capable
of producing the observed magnitude of wall thickening
\cite{maron2013}. Despite a heterogeneous clinical
spectrum -- many carriers remain asymptomatic for decades -- HCM is
recognised as a leading cause of sudden cardiac death (SCD) in
young people and competitive athletes, a risk that may manifest
without prior warning \cite{semsarian2015}. The combination of a
common genetic cause with a deceptively silent natural history makes
accurate variant-level genotype interpretation central to clinical
management.

The genetic architecture of HCM is dominated by mutations in genes
encoding components of the cardiac sarcomere, the fundamental
contractile unit of the cardiomyocyte \cite{seidman2001}. Roughly
60\% of genotype-positive HCM cases carry pathogenic or likely
pathogenic variants in sarcomeric protein-coding genes, with
\textit{MYH7} (encoding the $\beta$-myosin heavy chain) and
\textit{MYBPC3} (encoding cardiac myosin-binding protein C)
together accounting for over half of these cases \cite{seidman2001}.
The remaining disease burden is distributed across seven further
sarcomeric genes: \textit{TNNT2}, \textit{TNNI3}, \textit{TPM1},
\textit{ACTC1}, \textit{MYL2}, \textit{MYL3}, and \textit{TNNC1}.
These nine genes constitute the panel examined in this work.
Crucially, genotype--phenotype correlation in HCM is imperfect;
the same missense variant can produce widely differing clinical
presentations across carriers, complicating direct gene-level risk
inference and underscoring the need for variant-level pathogenicity
scoring.

\subsection{The VUS Bottleneck in HCM}

The clinical interpretation bottleneck in HCM lies not with frankly
pathogenic variants -- those supported by consistent functional
evidence and familial segregation data -- but with the large fraction
of observed missense substitutions that resist confident
classification. Under the ACMG/AMP framework \cite{richards2015},
variants that lack sufficient evidence to be placed in the
pathogenic or benign categories are designated Variants of Uncertain
Significance (VUS). In HCM panel testing, VUS findings are
particularly prevalent: the majority of detected missense
substitutions in sarcomeric genes cannot be definitively classified
with current evidence, leaving at-risk family members without
actionable genetic information, blocking predictive cascade
screening, and placing clinical teams in the position of managing
risk for an unresolved diagnosis.

The scale of this problem is not abstract. In the consolidated
multi-database variant set assembled for this study -- spanning 8,137
missense observations across the nine canonical HCM sarcomeric
genes -- more than 4,500 variants carry a VUS designation in ClinVar
\cite{landrum2016clinvar}, representing a clinically unresolved
majority that exceeds the total number of confidently labelled
variants by a factor greater than two. No existing computational
tool has been specifically designed to triage this backlog under
rigorous, leakage-free evaluation conditions. Resolving it is the
central motivation of this paper.

\subsection{Problem Statement and Approach}

\subsubsection{Label Circularity in Existing Predictors}

A fundamental flaw in many existing pathogenicity predictors is
\textit{label circularity} -- a form of data leakage in which model
inputs are co-derived from the same clinical workflows that generate
the training labels, causing the model to reconstruct the labelling
protocol rather than learn the underlying molecular biology. We
address this problem explicitly and at two levels in our design,
as detailed in Section~\ref{sec:related} and Section~\ref{sec:data}.

\subsubsection{Our Approach}

To address these limitations, we present a dual-model zero-leakage
pathogenicity prediction framework trained on 1,954 high-confidence
labelled missense variants across the nine sarcomeric HCM genes.
Variant labels are drawn exclusively from multi-submitter ClinVar
records \cite{landrum2016clinvar, landrum2018clinvar}; population
frequency data from the Genome Aggregation Database (gnomAD)
\cite{karczewski2020} supplement the benign class with
high-frequency population variants, correcting for the well-known
ascertainment bias in clinical databases toward pathogenic reports.
Structural and biochemical features are derived from UniProt
Swiss-Prot annotations \cite{uniprot2023}, yielding 51 engineered
descriptors per variant that capture domain localisation, secondary
structure context, binding site proximity, and physicochemical
substitution cost. All fields directly derivable from the ACMG/AMP
evidence criteria -- population frequency scores, clinical
significance annotations, review status, and prior disease
assertions -- are excluded from the feature set prior to training.

We train two complementary models: a Baseline Random Forest (RF)
operating on the 51 hand-crafted HCM-specific structural features,
and a Two-Tower Hybrid Neural Network that augments these tabular
features with differential evolutionary embeddings extracted from a
frozen 650-million-parameter ESM-2 protein language model
\cite{lin2023esm2}. Both models are evaluated under
\textit{Leave-One-Gene-Out} (LOGO) cross-validation, which withholds
all variants from a single gene at test time, preventing homology
contamination across the nine gene families. The resulting framework
is benchmarked against five established external tools spanning
evolutionary, structural, and ensemble paradigms.


% ==============================================================
% SECTION II  --  RELATED WORK
% (Teammate contribution  --  verbatim from Part1.tex)
% ==============================================================
\section{Related Work}
\label{sec:related}

The computational prediction of missense variant pathogenicity has
evolved substantially from early conservation-based heuristics to
large-scale deep learning frameworks, yet a persistent gap remains
between proteome-wide generality and the disease-specific precision
required for clinical variant interpretation in targeted gene panels.
We survey the three classes of prior work most directly relevant to
our contribution.

\subsection{Proteome-Wide Pathogenicity Prediction}

The dominant paradigm in computational variant interpretation has
been the development of tools trained and evaluated across the full
human proteome, exploiting signals -- evolutionary constraint,
structural context, allele frequency -- that are broadly available
regardless of the specific disease or gene under study.

Frazer \textit{et al.} introduced EVE (Evolutionary model of Variant
Effect) \cite{frazer2021eve}, a fully unsupervised variational
autoencoder trained on deep multiple sequence alignments drawn from
across the tree of life. By modelling the distribution of naturally
occurring sequence variation in each protein family, EVE infers
which positions and substitutions are evolutionarily tolerated and
which are constrained -- producing pathogenicity scores without
exposure to any clinical label. Applied to over 36 million variants
across 3,219 disease-associated human genes, EVE demonstrated that
evolutionary co-variation alone carries substantial pathogenicity
signal. However, because the model is blind to the disease-specific
structural context of sarcomeric proteins -- the precise domain
boundaries, regulatory interfaces, and force-transmission
architecture that define HCM pathomechanics -- its inferences are
necessarily generic.

Cheng \textit{et al.} presented AlphaMissense
\cite{cheng2023alphamissense}, an adaptation of the AlphaFold2
structure prediction framework fine-tuned on human and primate
population allele frequency data. By combining structural context
derived from AlphaFold's learned protein representations with
evolutionary information encoded in population-level variant
frequencies, AlphaMissense achieves proteome-wide coverage and
confidently classifies 89\% of all possible human missense
substitutions as likely pathogenic or likely benign. While
representing the current state of the art for proteome-wide
prediction, the model is trained on signals averaged across the
entire proteome and has not been adapted to the disease-specific
distribution of sarcomeric HCM variants; nor does its published
evaluation employ gene-held-out cross-validation.

Ioannidis \textit{et al.} proposed REVEL \cite{ioannidis2016revel},
an ensemble meta-predictor combining scores from thirteen individual
tools -- including SIFT, PolyPhen-2, MutPred, VEST, GERP++, and
phyloP -- using a random forest trained on ClinVar-derived labels.
Pre-computed REVEL scores are now standard components of clinical
bioinformatics pipelines, including dbNSFP and Ensembl VEP. REVEL's
design, however, embeds a structural dependence on ClinVar label
density: for genes that are heavily annotated in ClinVar, the
meta-predictor benefits from a rich prior that does not represent the
true challenge of unseen variant interpretation in less-studied
sarcomeric loci.

Li \textit{et al.} developed MetaRNN \cite{li2022metarnn}, a
recurrent neural network with gated recurrent units that integrates
28 functional annotation scores together with a local sequence
context window of $\pm$1 codon. Notably, MetaRNN explicitly avoids
incorporating scores derived from ClinVar-trained tools -- such as
MPC and ClinPred -- to circumvent what the authors term ``Type~I
circularity'': the inflation of performance metrics when training
and test variants overlap at the level of individual ClinVar
submissions. This represents a meaningful methodological step
forward. However, MetaRNN does not enforce gene-level isolation
between training and test splits, leaving the protein-level
homology contamination identified by Grimm \textit{et al.}
\cite{grimm2015circularity} -- ``Type~II circularity'' -- unaddressed.

While powerful at scale, these four proteome-wide tools share a
common limitation: they are trained and evaluated without controls
for gene-level homology leakage. For the specific challenge of
sarcomeric HCM variant interpretation, where domain localisation
and regulatory interface proximity carry outsized predictive weight,
this proteome-wide generality is a structural disadvantage rather
than a strength.

\subsection{Disease-Specific Tools for Inherited Cardiac Conditions}

The most direct predecessor to our work is CardioBoost
\cite{zhang2021cardioboost}, a gradient-boosted ensemble classifier
developed specifically for rare missense variants in genes associated
with inherited cardiomyopathies and arrhythmias. Trained exclusively
on variants with gnomAD allele frequency below 0.1\% -- a rarity
filter that substantially reduces benign contamination -- CardioBoost
demonstrated that disease-specific feature engineering and
domain-constrained training data yield meaningful improvements over
proteome-wide baselines, reporting an AUPRC of 91\% for
cardiomyopathy variants on its internal test set. The tool is
publicly available through the CardioClassifier platform and covers
the primary sarcomeric HCM genes, though it does not include
\textit{TNNC1} in its precomputed score set. Importantly,
CardioBoost was developed and internally validated using random
variant-level train-test splits rather than a gene-held-out
strategy -- a distinction that, as we demonstrate in
Section~\ref{sec:results}, has material implications for the
reliability of reported performance estimates.

\subsection{Data Resources}

The construction of a high-quality, leakage-free training dataset
for HCM variant pathogenicity requires the integration of three
complementary public resources: ClinVar \cite{landrum2016clinvar,
landrum2018clinvar} for clinical variant labels, UniProt Swiss-Prot
\cite{uniprot2023} for sequence and structural annotations, and the
Genome Aggregation Database (gnomAD) \cite{karczewski2020} for
population allele frequency data. The provenance, access method, and
exact role of each resource in our data pipeline are described in
detail in Section~\ref{sec:data}.

\subsection{Machine Learning for Variant Pathogenicity and the
Leakage Problem}

\subsubsection{The Case for Disease-Specific Machine Learning}

The comparative performance of CardioBoost against proteome-wide
baselines \cite{zhang2021cardioboost} established an important
empirical principle: when training data, feature engineering, and
model architecture are aligned with the specific disease context,
classification performance on targeted gene panels consistently
exceeds that of models trained on the full proteome. This advantage
arises because disease-relevant structural features -- sarcomere
domain membership, proximity to actin-binding or regulatory
interfaces, conservation within a cardiac-specific evolutionary
context -- are signals that a proteome-wide model can only approximate
through averaged cross-gene statistics, but that a disease-specific
model can encode directly.

The advance of protein language models (PLMs) adds a complementary
dimension to this argument. Lin \textit{et al.} demonstrated that
ESM-2 \cite{lin2023esm2}, a transformer trained on hundreds of
millions of protein sequences via masked language modelling, encodes
structural and evolutionary constraints with sufficient fidelity to
support atomic-level structure prediction from primary sequence
alone -- without explicit multiple sequence alignment at inference
time. These learned representations provide a continuous latent space
in which the evolutionary fitness landscape of a protein family is
implicitly encoded. A mutation-induced perturbation in this space,
derived from a frozen ESM-2 model, therefore carries interpretable
evolutionary signal without introducing any dependence on clinical
labels -- a property we exploit directly in the Two-Tower model
architecture described in Section~\ref{sec:data}.

\subsubsection{The Data Leakage Problem in Clinical Genomics}

A critical and frequently underaddressed methodological issue in the
development of variant pathogenicity predictors is data leakage
\cite{kapoor2023leakage}. In the clinical genomics setting, leakage
arises in two distinct forms. Grimm \textit{et al.}
\cite{grimm2015circularity} provided a foundational analysis of
both, identifying them as variant-level circularity (Type~I), in
which the same variants appear in training and evaluation sets, and
protein-level circularity (Type~II), in which variants from the same
protein populate both partitions even if the specific variants
differ. Grimm \textit{et al.} showed that failure to control for
either type leads to systematically inflated performance estimates
that do not reflect a model's ability to generalise to genuinely
novel variants \cite{grimm2015circularity}.

Feature-level leakage presents an orthogonal but equally serious
challenge: if model inputs include population allele frequency
annotations, prior clinical assertion counts, or ACMG evidence
tags -- all co-derived from the same expert workflow that generates
the training labels -- the model learns to recover the labelling
process rather than the underlying molecular biology. We address
both dimensions simultaneously: all ACMG-derivable fields are
excluded from the feature set, and model evaluation is conducted
under LOGO cross-validation which enforces complete gene-level
separation between training and test data.


% ==============================================================
% SECTION III  --  DATASET CONSTRUCTION AND FEATURE ENGINEERING
% (III-A, III-B verbatim from Part1.tex; III-C onward new)
% ==============================================================

% ==============================================================
% SECTION III  --  PROPOSED METHODOLOGY
% ==============================================================
\section{Proposed Methodology}
\label{sec:methodology}

Our methodology is designed to process multiple biological modalities while strictly avoiding the data leakage mechanisms that confound existing predictors. The architecture of our approach consists of robust data curation, zero-leakage feature extraction, a dual-model predictive framework, and a mechanistic interpretability layer.

\subsection{Data Curation and Problem Formulation}

To build a high-quality, zero-leakage training corpus, we established a rigorous data collection pipeline that spans clinical and functional annotations.

\vspace{2mm}\noindent\textbf{Data Sources.} 

The construction of the training dataset integrates three
complementary public resources, each contributing a distinct and
non-overlapping layer of evidence to the final variant
representation.

\textbf{ClinVar} \cite{landrum2016clinvar, landrum2018clinvar} is
the primary public archive for clinical interpretations of human
genetic variants, aggregating submissions from clinical testing
laboratories, research groups, and expert panels worldwide. Variants
in ClinVar are classified under the five-tier ACMG/AMP scheme, and
a review status field reflects the number and concordance of
independent submissions. For this study, only variants carrying a
two-star or higher review status -- indicating multi-submitter
consensus -- are admitted to the labelled training set. ClinVar also
contains the VUS records that constitute the applied target of this
work: the 4,500+ variants in our nine-gene panel whose clinical
resolution is the direct output objective of the deployed model.

\textbf{UniProt Swiss-Prot} \cite{uniprot2023} provides manually
curated annotations for each of the nine sarcomeric proteins in our
panel, including canonical amino acid sequences, domain boundary
coordinates, secondary structure element assignments (helix, strand,
coil), intrinsically disordered regions, binding site residues,
post-translational modification (PTM) sites, and active site flags.
UniProt and ClinVar data were accessed simultaneously via the
\textbf{EBI Proteins API}, a unified programmatic interface enabling
joint querying of both resources within a single data collection
pipeline.

\textbf{gnomAD} \cite{karczewski2020} aggregates whole-exome and
whole-genome sequencing data from approximately 730,000 individuals
across diverse ancestral backgrounds, providing allele frequency
estimates for the vast majority of observed human variants. In our
pipeline, variants with gnomAD allele frequency $> 0.1\%$ serve as
high-confidence benign proxies, addressing the well-documented
ascertainment bias in ClinVar toward pathogenic submissions and
ensuring that the benign class in the training set is credibly
populated rather than artificially sparse.

\vspace{2mm}\noindent\textbf{Dataset Statistics and Cohort Description.} 

The labeled dataset comprises \textbf{1,954 missense variants}
across nine sarcomeric genes: \textbf{1,441 pathogenic} (73.7\%)
and \textbf{513 benign} (26.3\%). The broader inference cohort spans
\textbf{8,137 variants} across the same genes, including
\textbf{4,524 VUS} targeted for probabilistic re-stratification.
Table~\ref{tab:cohort} summarises the per-gene composition.

\begin{table}[h]
\centering
\caption{Gene-wise Variant Counts in the Labeled Training Cohort}
\label{tab:cohort}
\begin{tabular}{lllr}
\toprule
\textbf{Gene} & \textbf{Protein} & \textbf{UniProt} & $N_{\mathrm{labeled}}$ \\
\midrule
\textit{MYH7}   & $\beta$-MHC       & P12883 & 986 \\
\textit{MYBPC3} & cMyBP-C          & Q14896 & 429 \\
\textit{TNNT2}  & Troponin T2      & P45379 & 130 \\
\textit{TNNI3}  & Troponin I3      & P19429 & 118 \\
\textit{TPM1}   & $\alpha$-Tropomyosin & P09493 & 100 \\
\textit{ACTC1}  & Cardiac Actin    & P68032 &  64 \\
\textit{MYL2}   & Regulatory MLC   & P10916 &  54 \\
\textit{MYL3}   & Essential MLC    & P08590 &  37 \\
\textit{TNNC1}  & Troponin C1      & P63316 &  36 \\
\midrule
\textbf{Total}  &                  &        & \textbf{1,954} \\
\bottomrule
\end{tabular}
\end{table}


\subsection{Feature Engineering and Embeddings}

To capture both the biophysical properties of the mutated residue and the broader evolutionary constraint on the protein, we engineered two distinct feature sets.

\vspace{2mm}\noindent\textbf{Zero-Leakage Protocol.} 

To eliminate feature-level label circularity, the following
columns are removed from the feature matrix before model training:
\texttt{pop\_freq}, \texttt{clin\_sig}, \texttt{review\_status},
\texttt{disease}, \texttt{sources}, and \texttt{genomic\_loc}.
Each field either directly encodes an ACMG/AMP classification criterion
or serves as a strong proxy for the clinical adjudication workflow that
generated the training labels. Their retention would allow a model to
reconstruct database consensus rather than learn intrinsic molecular
pathology. The remaining numeric and boolean features constitute the
51-dimensional tabular input to both classifiers.

\vspace{2mm}\noindent\textbf{The 51 Biophysical Descriptors.} 

The tabular feature vector $\xtab_i \in \mathbb{R}^{51}$ is derived
from UniProt structural annotations and encodes the local
biophysical microenvironment of each substitution across five
semantic groups:

\begin{itemize}
  \item \textbf{Physicochemical substitution cost.} The Grantham
    distance \cite{ng2003sift} between the wild-type and mutant amino
    acids captures multi-dimensional biochemical dissimilarity
    (composition $c$, polarity $p$, molecular volume $v$):

  \item \textbf{Charge and size perturbations.} Wild-type and mutant
    residue charge classes ($\in \{-1, 0, +1\}$) and size classes
    ($\in \{1,2,3\}$) are encoded, together with their signed deltas
    $\Delta\mathrm{charge}$ and $\Delta\mathrm{size}$.
  \item \textbf{Domain-level context.} Boolean flags indicate whether
    the mutated position resides within a motor domain, coiled-coil
    segment, regulatory zone, or other annotated structural region of
    the respective sarcomeric protein.
  \item \textbf{Critical site markers.} Separate binary features encode
    proximity to ligand-binding sites, metal-binding sites, active
    sites, phosphorylation / PTM sites, and secondary structure
    contexts (alpha-helix, beta-strand, coil).
  \item \textbf{Relative position.} The normalised residue position
    $\rho_i = \mathrm{pos}_i / L_{\mathrm{protein}}$ encodes depth
    within the polypeptide chain.
\end{itemize}

Missing tabular values (resulting from partial UniProt annotation
coverage) are imputed to zero, a conservative choice consistent with
the reproducible implementation.

\vspace{2mm}\noindent\textbf{Differential ESM-2 Embeddings.} 

A core input to the Two-Tower model is the \textit{differential
evolutionary embedding} $\Deltaemb_i$ derived from a frozen instance
of ESM-2 \cite{lin2023esm2} (variant \texttt{esm2\_t33\_650M\_UR50D},
$33$ Transformer layers, $650\,\mathrm{M}$ parameters). For each
missense variant~$i$, the wild-type sequence context $s_i^{\mathrm{WT}}$
and mutant sequence context $s_i^{\mathrm{MUT}}$ are independently
encoded by the frozen encoder $f_\theta$. Mean pooling over the
sequence length dimension yields 1,280-dimensional protein-level
representations:
\begin{equation}
  \mathbf{e}_i^{\mathrm{WT}}  = \mathrm{MeanPool}\bigl(f_\theta(s_i^{\mathrm{WT}})\bigr),\quad
  \mathbf{e}_i^{\mathrm{MUT}} = \mathrm{MeanPool}\bigl(f_\theta(s_i^{\mathrm{MUT}})\bigr).
\end{equation}
The mutation-induced perturbation vector is then
\begin{equation}
  \Deltaemb_i = \mathbf{e}_i^{\mathrm{MUT}} - \mathbf{e}_i^{\mathrm{WT}}
  \;\in\;\mathbb{R}^{1280}.
  \label{eq:delta_emb}
\end{equation}
Subtracting the wild-type embedding attenuates protein-family
background signal and isolates the mutation-local shift in
evolutionary latent space \cite{frazer2021eve}. Because ESM-2
parameters remain strictly frozen throughout training, the embeddings
carry no statistical dependence on the clinical labels used to
supervise the downstream classifier, preserving the zero-leakage
guarantee. Pre-computed embeddings for all 1,954 labelled variants
are stored in \texttt{data/esm2\_delta\_embeddings.npy}
($1954 \times 1280$ float32 array).

\subsection{Two-Tower Hybrid Neural Network}

Our primary predictive engine is a custom dual-input architecture designed to fuse tabular biophysical descriptors with high-dimensional language model embeddings. By processing these two distinct modalities through dedicated projection pathways before integrating them, the network preserves the unique structural constraints of the sarcomere while simultaneously leveraging the rich evolutionary context encoded by the protein language model. This late-fusion strategy prevents early information bottlenecking and allows both domains to contribute optimally to the final pathogenicity score.

\vspace{2mm}\noindent\textbf{Architecture Overview.} 

The \textsc{HybridHCMModel} is a custom two-input neural network
that processes the tabular biophysical descriptor vector and the
ESM-2 differential embedding through two independent linear
projections (``towers''), concatenates the resulting representations,
and passes the fused vector through a shared prediction head.
All layers use hidden width $h = 32$ and apply dropout at rate
$p_\mathrm{drop} = 0.3$ after each activation.

\vspace{2mm}\noindent\textbf{Tabular Feature Tower.} 

For variant $i$, the tabular input $\xtab_i \in \mathbb{R}^{d_\mathrm{tab}}$
($d_\mathrm{tab}=51$) is transformed by
\begin{equation}
  \ztab_i = \mathrm{Dropout}\!\left(
      \mathrm{BN}\!\left(
        \mathrm{ReLU}\!\left(W_\mathrm{tab}\,\xtab_i + \mathbf{b}_\mathrm{tab}\right)
      \right)
    \right),
  \label{eq:tab_tower}
\end{equation}
where $W_\mathrm{tab}\in\mathbb{R}^{h\times d_\mathrm{tab}}$,
$\mathbf{b}_\mathrm{tab}\in\mathbb{R}^h$, and BN denotes batch
normalisation with learnable scale and shift parameters.

\vspace{2mm}\noindent\textbf{ESM-2 Embedding Tower.} 

The differential embedding $\xesm_i = \Deltaemb_i \in \mathbb{R}^{1280}$
(Eq.~\ref{eq:delta_emb}) passes through an identically structured
projection:
\begin{equation}
  \zesm_i = \mathrm{Dropout}\!\left(
      \mathrm{BN}\!\left(
        \mathrm{ReLU}\!\left(W_\mathrm{esm}\,\xesm_i + \mathbf{b}_\mathrm{esm}\right)
      \right)
    \right),
  \label{eq:esm_tower}
\end{equation}
where $W_\mathrm{esm}\in\mathbb{R}^{h\times 1280}$ and
$\mathbf{b}_\mathrm{esm}\in\mathbb{R}^h$.

\vspace{2mm}\noindent\textbf{Fusion and Output Head.} 

The two tower representations are concatenated along the feature
dimension to form the fused representation
$\zfus_i = [\ztab_i\,;\,\zesm_i] \in \mathbb{R}^{2h}$.
The prediction head applies a further linear projection with
batch normalisation, dropout, and a final scalar projection through
a sigmoid activation:
\begin{align}
  \mathbf{h}_i        &= \mathrm{Dropout}\!\left(\mathrm{BN}\!\left(\mathrm{ReLU}(W_f\,\zfus_i + \mathbf{b}_f)\right)\right),
                         \quad W_f\in\mathbb{R}^{h\times 2h}, \label{eq:fusion}\\
  \hat{p}_i^{\,\mathrm{unc}}  &= \sigma\!\left(w_o^\top\mathbf{h}_i + b_o\right),
                         \quad w_o\in\mathbb{R}^h, \label{eq:output}
\end{align}
where $\sigma(x)=(1+e^{-x})^{-1}$ is the logistic sigmoid. The
uncalibrated output $\hat{p}_i^{\,\mathrm{unc}}\in(0,1)$ is an
initial pathogenicity probability estimate that is subsequently
refined by Platt scaling (Section~\ref{sec:platt}).

\vspace{2mm}\noindent\textbf{Class-Balanced Binary Cross-Entropy Loss.} 

The labeled training set contains 1,441 pathogenic and 513 benign
variants, yielding a 2.8:1 positive-to-negative ratio. To prevent
the network from deferring to the majority class, per-sample
inverse-frequency weights are computed from the training split
$\{y_i\}_{i\in\mathcal{D}_\mathrm{train}}$:
\begin{equation}
  w_i = \begin{cases}
    \dfrac{N}{2N_1} & \text{if } y_i = 1,\\[6pt]
    \dfrac{N}{2N_0} & \text{if } y_i = 0,
  \end{cases}
  \label{eq:weights}
\end{equation}
where $N = N_1 + N_0$ is the total training-set size. The weighted
Binary Cross-Entropy (BCE) objective is
\begin{equation}
  \mathcal{L}_\mathrm{BCE} = -\frac{1}{N}\sum_{i=1}^{N}
    w_i\!\left[y_i \log \hat{p}_i^{\,\mathrm{unc}} +
    (1-y_i)\log(1-\hat{p}_i^{\,\mathrm{unc}})\right].
  \label{eq:bce}
\end{equation}
The network is optimised with AdamW \cite{loshchilov2019adamw}
(learning rate $\eta=0.003$, weight-decay $\lambda=10^{-4}$) for
$T=40$ epochs with full-batch gradient computation.

\subsection{Baseline Random Forest Classifier}

To establish a strong traditional baseline against our neural architecture and to evaluate the necessity of cross-modal fusion, we train an ensemble decision tree classifier.

\vspace{2mm}\noindent\textbf{Ensemble Formulation.} 

To establish a strong traditional baseline, we train a Random Forest
classifier \cite{breiman2001rf}. Unlike the Two-Tower model, which
processes each modality through a dedicated projection, the Random
Forest receives the \emph{concatenated} representation
\begin{equation}
  \mathbf{x}_i^{\mathrm{RF}} = \bigl[\xtab_i\,;\,\xesm_i\bigr]
    \in \mathbb{R}^{d_\mathrm{tab}+1280},
  \label{eq:rf_input}
\end{equation}
treating the combined feature space as a flat input to an ensemble
of $B=200$ Classification and Regression Trees (CARTs).

\vspace{2mm}\noindent\textbf{Node Split Criterion: Gini Impurity.} 

At each internal node $t$ during tree induction, the optimal split
is selected by maximising the reduction in \emph{Gini impurity}:
\begin{equation}
  I_G(t) = 1 - \sum_{k\in\{0,1\}} \hat{p}(k|t)^2,
  \label{eq:gini}
\end{equation}
where $\hat{p}(k|t)$ is the empirical proportion of class-$k$
instances at node $t$. The split that produces the largest weighted
reduction
\begin{equation}
  \Delta I_G = I_G(t) - \frac{n_L}{n}I_G(t_L) - \frac{n_R}{n}I_G(t_R)
\end{equation}
is selected, where $t_L, t_R$ are the left and right child nodes
containing $n_L$ and $n_R$ samples, respectively. Tree depth is
capped at $\mathrm{max\_depth}=10$ to control variance.

\vspace{2mm}\noindent\textbf{Class-Balanced Bootstrap Sampling.} 

To address label imbalance within each bootstrap sample, class-weight
balanced sampling is used: effective sample weights are inversely
proportional to class frequencies in each bootstrap draw. The
aggregated probability estimate across the ensemble is
\begin{equation}
  \hat{p}_i^{\,\mathrm{RF}} = \frac{1}{B}\sum_{b=1}^{B} T_b\!\left(\mathbf{x}_i^{\mathrm{RF}}\right),
  \label{eq:rf_pred}
\end{equation}
where $T_b(\cdot)$ returns the class-1 probability from tree $b$.

\vspace{2mm}\noindent\textbf{Purpose of the Baseline.} 

The Baseline RF answers a specific scientific question: can a
traditional ensemble model fully exploit the high-dimensional ESM-2
evolutionary signal through simple concatenation and axis-aligned
splits, without specialised cross-modal fusion? Comparing its
performance against the Two-Tower model in identical LOGO conditions
isolates the incremental value of learned multi-modal integration
over bag-of-features treatment of the same information.


% ==============================================================
% SECTION VI  --  EXPLAINABLE AI VIA IN SILICO MUTAGENESIS
% ==============================================================
\subsection{Explainability via \textit{In Silico} Mutagenesis}

Pointwise prediction without mechanistic interpretability is insufficient for clinical deployment. To expose the biochemical logic encoded by each model, we deploy an \textit{in silico} mutagenesis framework.

\vspace{2mm}\noindent\textbf{ISM Landscape Generation.} 

Pointwise prediction without mechanistic interpretability is
insufficient for clinical deployment. To expose the biochemical
logic encoded by each model, we generate \textit{in silico}
mutagenesis (ISM) landscapes: comprehensive maps of predicted
pathogenicity across all possible single-amino-acid substitutions
in a target sequence segment.

For a gene-representative sequence window of length $L$ and
canonical amino-acid alphabet $\mathcal{A}$ ($|\mathcal{A}|=20$),
we systematically construct mutant contexts
$s^{(j,a)}$ for each position $j\in\{1,\ldots,L\}$ and alternative
amino acid $a\in\mathcal{A}$. For each such mutant context, we:
\begin{enumerate}
  \item Compute the differential ESM-2 embedding
    $\Deltaemb^{(j,a)} = f_\theta(s^{(j,a)})_j - f_\theta(s^{\mathrm{WT}})_j \in \mathbb{R}^{1280}$.
  \item Update the tabular descriptor vector to reflect the physical
    properties of substitution $(a_j \to a)$: Grantham distance,
    charge delta, size delta, and local chemical context.
  \item Forward-pass the updated $(\mathbf{x}^{\mathrm{tab},(j,a)},\,\Deltaemb^{(j,a)})$
    through the calibrated Two-Tower predictor (or Baseline RF).
\end{enumerate}
This yields the ISM matrix $M \in [0,1]^{20 \times L}$:
\begin{equation}
  M_{a,j} = \hat{p}^{\,\mathrm{cal}}\!\left(\mathbf{x}^{\mathrm{tab},(j,a)},\,
      \Deltaemb^{(j,a)}\right),
  \label{eq:ism_matrix}
\end{equation}
with the wild-type diagonal $M_{a_j, j}$ serving as the positional
baseline.

\vspace{2mm}\noindent\textbf{Heatmap Visualisation and Biomechanical Mapping.} 

The ISM matrix is visualised as a \textit{coolwarm} heatmap
(rows: 20 amino acids; columns: residue positions), where warm
colours indicate high predicted pathogenicity and cool colours
indicate predicted benignity. Vertical hot-spots -- residue positions
where nearly every substitution is predicted damaging -- identify
structurally intolerant positions. Horizontal stripes of warmth
reveal substitution classes (e.g. charge-reversing or steric-clashing)
that are universally destabilising.

Crucially, ISM landscapes are generated from a model trained
\emph{without} any explicit 3-D structural input. The emergence of
coherent spatial patterns that align with known sarcomeric structural
biology therefore constitutes evidence that the combined
ESM-2~+~biophysical representation has learned meaningful
sequence-structure-function relationships. Mapping ISM hot-spots
against UniProt domain annotations allows us to correlate high
predicted pathogenicity with known functional interfaces (e.g.\ the
myosin converter domain, actin-binding cleft, or troponin regulatory
interface), providing biologically interpretable hypotheses for
prioritisation of functional validation experiments.


% ==============================================================
% SECTION IV  --  EXPERIMENTAL DESIGN AND EVALUATION
% ==============================================================
\section{Experimental Design and Evaluation}
\label{sec:experimental_setup}

To rigorously assess model performance and avoid data leakage, we utilize the following evaluation protocols and metrics.

\vspace{2mm}\noindent\textbf{Leave-One-Gene-Out Validation Protocol.} 

To prevent Type-II homology leakage \cite{grimm2015circularity},
model evaluation is conducted under Leave-One-Gene-Out (LOGO)
cross-validation. For the nine-gene target set
$\mathcal{G} = \{\textit{MYH7},\,\textit{MYBPC3},\,\textit{TNNT2},\dots\}$
and held-out gene $g^* \in \mathcal{G}$:
\begin{align}
  \mathcal{D}_\mathrm{train} &= \bigl\{(\mathbf{x}_i,\,y_i) : g_i \neq g^*\bigr\}, \\
  \mathcal{D}_\mathrm{test}  &= \bigl\{(\mathbf{x}_i,\,y_i) : g_i  =  g^*\bigr\}.
\end{align}
This procedure is repeated nine times, each time holding out a
different gene, so that every variant in the dataset is predicted
exactly once by a model that has never seen another variant from the
same protein family.

\vspace{2mm}\noindent\textbf{Bootstrap Confidence Intervals.} 

Point estimates of the Area Under the Precision-Recall Curve (AUPRC)
are supplemented by nonparametric 95\% bootstrap confidence intervals.
For $B = 100$ resampling rounds on the LOGO test set
$\{(y_i, \phat_i)\}_{i=1}^n$:
\begin{equation}
  \overline{\mathrm{AUPRC}} = \frac{1}{|\mathcal{B}|}\sum_{b\in\mathcal{B}}\mathrm{AUPRC}^{(b)}, \quad
  \mathrm{CI}_{95\%} = \bigl[Q_{2.5},\,Q_{97.5}\bigr].
\end{equation}
Bootstrapped samples containing fewer than two distinct classes are
discarded, consistent with the implementation safeguard.

\vspace{2mm}\noindent\textbf{Probability Calibration via Platt Scaling.}\label{sec:platt} 

Neural sigmoid outputs may be systematically over- or under-confident,
particularly when trained with class-reweighted objectives. For
clinical decision support, predictions must represent calibrated
probabilities. We apply Platt scaling: a univariate logistic
regression $\mathcal{C}$ is fitted on the \emph{in-fold} training
predictions $\hat{p}_i^{\,\mathrm{unc}}$ (clipped to
$[10^{-7},1-10^{-7}]$ for numerical stability) to learn scalar
parameters $A,B \in \mathbb{R}$:
\begin{equation}
  \hat{p}_i^{\,\mathrm{cal}} = \mathcal{C}\!\left(\hat{p}_i^{\,\mathrm{unc}}\right)
    = \sigma\!\left(A\cdot\hat{p}_i^{\,\mathrm{unc}} + B\right).
  \label{eq:platt}
\end{equation}
Parameters are estimated by maximum likelihood with a large
regularisation constant ($C=999$) to prevent shrinkage.
Calibration quality is quantified by the Brier score,
\begin{equation}
  \mathrm{BS} = \frac{1}{n}\sum_{i=1}^{n}(\hat{p}_i^{\,\mathrm{cal}} - y_i)^2,
  \label{eq:brier}
\end{equation}
and the Expected Calibration Error (ECE), computed as the
sample-weighted mean absolute deviation between bin-wise mean
predicted probabilities and bin-wise observed positive fractions
across ten equal-width probability bins.


\section{Results}
\label{sec:results}

\subsection{Internal LOGO Validation: Two-Tower vs.\ Baseline RF}

Table~\ref{tab:logo} reports AUPRC, 95\% bootstrap confidence
intervals, Brier score, and Expected Calibration Error for both
models across all nine genes. 

\begin{table*}[t]
\centering
\caption{Leave-One-Gene-Out Cross-Validation Performance:
  Two-Tower Hybrid vs.\ Baseline Random Forest.
  AUPRC and its 95\% bootstrap CI are the primary discrimination
  metrics; Brier and ECE quantify probabilistic calibration.}
\label{tab:logo}
\setlength{\tabcolsep}{4.5pt}
\begin{tabular}{l cc cc cc cc}
\toprule
& \multicolumn{4}{c}{\textbf{Two-Tower Hybrid}} &
  \multicolumn{4}{c}{\textbf{Baseline Random Forest}} \\
\cmidrule(lr){2-5}\cmidrule(lr){6-9}
\textbf{Held-Out Gene} &
  \textbf{AUPRC} & \textbf{95\% CI} & \textbf{Brier} & \textbf{ECE} &
  \textbf{AUPRC} & \textbf{95\% CI} & \textbf{Brier} & \textbf{ECE} \\
\midrule
\textit{MYH7}   & 0.8636 & [0.835, 0.887] & 0.596 & 0.652 &
                  \textbf{0.9086} & [0.881, 0.931] & \textbf{0.167} & 0.183 \\
\textit{MYBPC3} & 0.6061 & [0.544, 0.666] & 0.355 & 0.344 &
                  \textbf{0.6412} & [0.575, 0.711] & \textbf{0.273} & \textbf{0.227} \\
\textit{TNNT2}  & \textbf{0.9206} & [0.869, 0.954] & \textbf{0.145} & \textbf{0.044} &
                  0.9205 & [0.852, 0.962] & 0.161 & 0.158 \\
\textit{TNNI3}  & \textbf{0.9420} & [0.871, 0.986] & \textbf{0.108} & \textbf{0.077} &
                  0.9323 & [0.879, 0.974] & 0.142 & 0.151 \\
\textit{TPM1}   & 0.9034 & [0.833, 0.974] & \textbf{0.124} & \textbf{0.050} &
                  \textbf{0.9649} & [0.929, 0.991] & 0.125 & 0.179 \\
\textit{ACTC1}  & 0.9400 & [0.880, 0.987] & \textbf{0.122} & \textbf{0.064} &
                  \textbf{0.9664} & [0.920, 0.993] & 0.136 & 0.172 \\
\textit{MYL2}   & \textbf{0.8742} & [0.740, 0.964] & 0.221 & 0.187 &
                  0.8593 & [0.754, 0.962] & \textbf{0.184} & \textbf{0.116} \\
\textit{MYL3}   & 0.8261 & [0.654, 0.933] & 0.266 & 0.271 &
                  \textbf{0.8460} & [0.728, 0.940] & \textbf{0.226} & \textbf{0.160} \\
\textit{TNNC1}  & 0.8324 & [0.667, 0.962] & 0.207 & 0.167 &
                  \textbf{0.9339} & [0.848, 0.995] & \textbf{0.156} & \textbf{0.127} \\
\midrule
\textbf{Mean}   & 0.8565 &  --  & 0.227 & 0.200 &
                  \textbf{0.8859} &  --  & \textbf{0.186} & \textbf{0.164} \\
\bottomrule
\end{tabular}
\end{table*}

The Baseline RF achieves a mean AUPRC of \textbf{0.886}, outperforming
the Two-Tower Hybrid (\textbf{0.857}) in seven of nine genes. The
Two-Tower model leads on \textit{TNNI3} (0.942 vs 0.932),
\textit{TNNT2} (0.921 vs 0.921, essentially tied), and \textit{MYL2}
(0.874 vs 0.859). Notably, the Two-Tower model achieves substantially
better calibration on most troponin and thin-filament genes
(TNNT2 ECE 0.044, TNNI3 ECE 0.077, TPM1 ECE 0.050), while the RF
maintains better calibration in myosin heavy chain and myosin-binding
protein C genes -- reflecting the complementary strengths of neural
calibration with Platt scaling versus tree-ensemble probability
smoothing.

\subsection{Reliability and Calibration Analysis}

Reliability diagrams for the \textit{TNNT2} LOGO hold-out cohort
are shown in Fig.~\ref{fig:cal_tt} (Two-Tower) and
Fig.~\ref{fig:cal_rf} (Baseline RF). The Platt-calibrated Two-Tower
model achieves a Brier score of \textbf{0.145} and ECE of
\textbf{0.044} on this held-out cohort, indicating that the predicted
probability bins closely track the empirical positive fraction. The
Baseline RF shows a comparable Brier score (0.161) but higher ECE
(0.158), suggesting mild overconfidence in intermediate probability
bins that Platt scaling addresses in the neural model.

\begin{figure}[h]
\centering
\includegraphics[width=\columnwidth]{calibration_plot_m2_Two-Tower_Hybrid.png}
\caption{Reliability diagram (TNNT2 LOGO hold-out): Platt-calibrated
Two-Tower Hybrid model. Brier$=0.145$, ECE$=0.044$.}
\label{fig:cal_tt}
\end{figure}

\begin{figure}[h]
\centering
\includegraphics[width=\columnwidth]{calibration_plot_m2_Baseline_RF.png}
\caption{Reliability diagram (TNNT2 LOGO hold-out): Baseline Random
Forest. Brier$=0.161$, ECE$=0.158$.}
\label{fig:cal_rf}
\end{figure}

\subsection{Ablation Study: Quantifying Modality Contributions}

To isolate the contributions of the biophysical and language-model
branches, we performed a modality ablation on the
\textit{TNNT2} held-out cohort by training separate models with one
input zeroed:

\begin{itemize}
  \item \textbf{Full fusion model:} point-estimate AUPRC \textbf{0.931}.
  \item \textbf{Tabular-only} (ESM branch zeroed):
    AUPRC \textbf{0.870}.
  \item \textbf{ESM-only} (tabular branch zeroed):
    AUPRC \textbf{0.940}.
\end{itemize}

The ESM-only ablation marginally outperforms the full fusion model
($+0.009$ AUPRC), while the tabular-only model lags by $0.061$
AUPRC. This finding demonstrates that the 650M-parameter evolutionary
language model encodes the vast majority of pathogenicity-relevant
signal within its latent representations. The strong tabular-only
performance (0.870) confirms that the 51 handcrafted structural
features carry substantial discriminative content; however, the
ESM representations appear to subsume and extend this information
into a higher-dimensional embedding space that simple concatenation
does not always improve upon.

\subsection{External Benchmarking: EVE and AlphaMissense}

Table~\ref{tab:bench_eve} compares our models against EVE and
AlphaMissense. All comparisons are evaluated on identical LOGO test
splits with matched variant subsets. Both our models outperform EVE
and AlphaMissense in seven of nine genes. Mean AUPRC gains of the
Baseline RF over EVE ($+0.095$) and AlphaMissense ($+0.086$) are
substantial and consistent.

\begin{table}[h]
\centering
\caption{AUPRC Comparison: Our Models vs.\ EVE and AlphaMissense.
  Best value per gene shown in \textbf{bold}.}
\label{tab:bench_eve}
\resizebox{\columnwidth}{!}{%
\begin{tabular}{lcccc}
\toprule
\textbf{Gene} & \textbf{Two-Tower} & \textbf{Baseline RF} &
  \textbf{EVE} & \textbf{AlphaMissense} \\
\midrule
\textit{ACTC1}  & 0.955 & \textbf{0.966} & 0.885 & 0.874 \\
\textit{MYBPC3} & \textbf{0.667} & 0.641 & 0.636 & 0.643 \\
\textit{MYH7}   & 0.886 & \textbf{0.909} & 0.873 & 0.880 \\
\textit{MYL2}   & 0.835 & \textbf{0.859} & 0.723 & 0.733 \\
\textit{MYL3}   & 0.806 & \textbf{0.846} & 0.624 & 0.616 \\
\textit{TNNC1}  & 0.693 & \textbf{0.934} & 0.811 & 0.840 \\
\textit{TNNI3}  & \textbf{0.938} & 0.932 & 0.887 & 0.894 \\
\textit{TNNT2}  & \textbf{0.931} & 0.921 & 0.812 & 0.828 \\
\textit{TPM1}   & 0.886 & \textbf{0.965} & 0.865 & 0.894 \\
\midrule
\textbf{Mean}   & 0.843 & \textbf{0.886} & 0.791 & 0.800 \\
\bottomrule
\end{tabular}%
}
\end{table}

\subsection{External Benchmarking: REVEL and MetaRNN}

Table~\ref{tab:bench_revel} presents the comparison with REVEL and
MetaRNN. Coverage is limited by dbNSFP annotation density: matched
variant counts range from $n=9$ (\textit{TNNT2}) to $n=74$
(\textit{MYH7}). Our Baseline RF exceeds both tools in mean AUPRC
(RF 0.886 vs REVEL 0.842 and MetaRNN 0.849), and leads in five of
nine genes. Two notable exceptions deserve attention: on
\textit{MYH7}, REVEL and MetaRNN achieve AUPRC $\approx 0.972$
versus our RF's 0.909; on \textit{TNNI3}, MetaRNN achieves 0.976
versus our RF's 0.932. Both genes are among the most densely
annotated in ClinVar, and REVEL/MetaRNN incorporate allele-frequency
evidence that is correlated with these labels. As discussed in
Section~\ref{sec:discussion}, this reflects ClinVar saturation
rather than superior molecular discrimination.

\begin{table}[h]
\centering
\caption{AUPRC Comparison: Our Models vs.\ REVEL and MetaRNN.
  $n$ denotes the number of variants matched to dbNSFP.
  Best value per gene shown in \textbf{bold}.}
\label{tab:bench_revel}
\resizebox{\columnwidth}{!}{%
\begin{tabular}{lccccc}
\toprule
\textbf{Gene} & \textbf{TT} & \textbf{Baseline RF} &
  \textbf{REVEL} & \textbf{MetaRNN} & $n$ \\
\midrule
\textit{ACTC1}  & 0.908 & \textbf{0.966} & 0.860 & 0.879 & 28 \\
\textit{MYBPC3} & 0.628 & \textbf{0.641} & 0.592 & 0.587 & 51 \\
\textit{MYH7}   & 0.871 & 0.909 & \textbf{0.972} & 0.972 & 74 \\
\textit{MYL2}   & 0.833 & \textbf{0.859} & 0.873 & 0.868 & 44 \\
\textit{MYL3}   & \textbf{0.815} & 0.846 & 0.882 & 0.754 & 28 \\
\textit{TNNC1}  & 0.805 & \textbf{0.934} & 0.749 & 0.728 & 30 \\
\textit{TNNI3}  & 0.910 & 0.932 & 0.959 & \textbf{0.976} & 61 \\
\textit{TNNT2}  & 0.911 & \textbf{0.921} & 0.755 & 0.938 &  9 \\
\textit{TPM1}   & 0.904 & \textbf{0.965} & 0.934 & 0.942 & 22 \\
\midrule
\textbf{Mean}   & 0.843 & \textbf{0.886} & 0.842 & 0.849 &  --  \\
\bottomrule
\end{tabular}%
}
\end{table}

\subsection{External Benchmarking: CardioBoost}

Table~\ref{tab:bench_cb} compares our models against CardioBoost,
the most directly competing disease-specific tool. CardioBoost scores
are unavailable for \textit{TNNC1} (zero matched variants), which is
excluded from the mean calculation. Our Baseline RF (mean AUPRC
0.883 over eight comparable genes) outperforms CardioBoost (0.876).
CardioBoost leads in two genes: \textit{MYBPC3} (0.776 vs 0.641 RF)
and \textit{MYH7} (0.939 vs 0.909 RF). Both advantages can be
attributed to the CardioBoost training strategy: \textit{MYBPC3} is
dominated by truncating variants for which CardioBoost was explicitly
engineered, and \textit{MYH7} has the densest ClinVar label coverage
of any panel gene, benefiting tools with implicit ClinVar-prior
alignment. Our models outperform CardioBoost in the remaining six
comparable genes.

\begin{table}[h]
\centering
\caption{AUPRC Comparison: Our Models vs.\ CardioBoost.
  N/A indicates no CardioBoost scores available for \textit{TNNC1}.
  Best value per gene shown in \textbf{bold}.}
\label{tab:bench_cb}
\resizebox{\columnwidth}{!}{%
\begin{tabular}{lccccc}
\toprule
\textbf{Gene} & \textbf{TT} & \textbf{Baseline RF} &
  \textbf{CardioBoost} & \textbf{Coverage} & $n$ \\
\midrule
\textit{ACTC1}  & 0.955 & \textbf{0.966} & 0.886 & 100.0\% &  64 \\
\textit{MYBPC3} & 0.667 & 0.641 & \textbf{0.776} &  95.1\% & 410 \\
\textit{MYH7}   & 0.886 & 0.909 & \textbf{0.939} &  99.1\% & 983 \\
\textit{MYL2}   & 0.835 & \textbf{0.859} & 0.763 &  98.1\% &  53 \\
\textit{MYL3}   & 0.806 & \textbf{0.846} & 0.832 &  97.3\% &  36 \\
\textit{TNNC1}  & 0.693 & \textbf{0.934} & N/A   &     --   &   0 \\
\textit{TNNI3}  & \textbf{0.938} & 0.932 & \textbf{0.952} &  97.5\% & 115 \\
\textit{TNNT2}  & \textbf{0.921} & 0.921 & 0.905 &  96.3\% & 129 \\
\textit{TPM1}   & 0.903 & \textbf{0.965} & \textbf{0.969} &  79.0\% &  79 \\
\midrule
\textbf{Mean}$^*$ & 0.851 & \textbf{0.883} & 0.876 & 95.0\% &  --  \\
\bottomrule
\multicolumn{6}{l}{\small $^*$Excluding \textit{TNNC1} for which CardioBoost has no coverage.}
\end{tabular}%
}
\end{table}

\subsection{Summary: Overall Benchmarking Performance}

Table~\ref{tab:bench_summary} consolidates the cross-tool comparison
into a single overview. The Baseline RF is the consistently
strongest performer across the full panel, with positive AUPRC
deltas against all five external tools. The Two-Tower Hybrid model
surpasses the proteome-wide tools and is competitive with
CardioBoost, while offering superior probabilistic calibration.

\begin{table}[h]
\centering
\caption{Overall External Benchmarking Summary. $\Delta$RF and
$\Delta$TT are mean-AUPRC advantages of our Baseline RF and
Two-Tower Hybrid relative to each external tool.}
\label{tab:bench_summary}
\begin{tabular}{lcccc}
\toprule
\textbf{Tool} & \textbf{Scope} & \textbf{Mean AUPRC} &
  $\boldsymbol{\Delta}$\textbf{RF} & $\boldsymbol{\Delta}$\textbf{TT} \\
\midrule
EVE             & Proteome-wide & 0.791 & \textbf{+0.095} & +0.052 \\
AlphaMissense   & Proteome-wide & 0.800 & \textbf{+0.086} & +0.043 \\
REVEL           & Proteome-wide & 0.842 & \textbf{+0.044} & +0.001 \\
MetaRNN         & Proteome-wide & 0.849 & \textbf{+0.037} & $-$0.006 \\
CardioBoost     & HCM-specific  & 0.876 & \textbf{+0.007} & $-$0.025 \\
\midrule
\textbf{Baseline RF}   & HCM-specific & \textbf{0.886} &  --  &  --  \\
\textbf{Two-Tower}     & HCM-specific & \textbf{0.857} &  --  &  --  \\
\bottomrule
\end{tabular}
\end{table}

\subsection{XAI Case Study: ISM Landscapes for MYH7 and TNNT2}

Fig.~\ref{fig:ism_myh7} and Fig.~\ref{fig:ism_tnnt2} display the
Two-Tower ISM landscapes for \textit{MYH7} and \textit{TNNT2},
respectively. In the \textit{MYH7} landscape, the model identifies
discrete residue positions where nearly every substitution is
predicted highly pathogenic (cell values $>0.90$). These hot-spot
columns co-localise with the myosin converter domain (UniProt
annotation positions 699--790 of P12883), a structurally constrained
mechanical socket that translates ATP hydrolysis into lever-arm
rotation. Substitutions at these positions are predicted pathogenic
irrespective of the specific alternative amino acid, consistent with
the functional intolerance of converter-domain geometry to
perturbation \cite{spudich2019}.

Horizontal rows of elevated risk in the \textit{MYH7} heatmap
correspond to charge-reversing residues (D, E $\to$ K, R)
and bulky hydrophobic substitutions at surface-exposed positions,
reflecting the model's implicit learning of
electrostatic and steric constraints in the actin-binding cleft.
The \textit{TNNT2} landscape shows a distinct hot-spot pattern
concentrated in the troponin T2 N-terminal regulatory region, which
mediates calcium-sensitive interaction with tropomyosin. Critically,
both structural interpretations emerge without explicit 3-D
coordinate supervision, confirming that the ESM-2 evolutionary
embedding has encoded meaningful sequence-structure-function
relationships.

\begin{figure}[h]
\centering
\includegraphics[width=\columnwidth]{ism_landscape_m2_TwoTower_MYH7.png}
\caption{\textit{In silico} mutagenesis (ISM) heatmap: Two-Tower
  Hybrid model, \textit{MYH7} (20 representative positions). Warm
  colours (red) indicate high predicted pathogenicity ($\hat{p}$);
  cool colours (blue) indicate predicted benignity. Vertical
  hot-spot bands correspond to the myosin converter domain.}
\label{fig:ism_myh7}
\end{figure}

\begin{figure}[h]
\centering
\includegraphics[width=\columnwidth]{ism_landscape_m2_TwoTower_TNNT2.png}
\caption{\textit{In silico} mutagenesis (ISM) heatmap: Two-Tower
  Hybrid model, \textit{TNNT2}. Hot-spot columns correspond to the
  N-terminal tropomyosin-binding regulatory region.}
\label{fig:ism_tnnt2}
\end{figure}

\subsection{High-Throughput VUS Prioritisation}

Applying the trained and calibrated Two-Tower model to 4,524
unresolved ClinVar VUS produces a rank-ordered list of candidates
stratified by predicted pathogenicity. Table~\ref{tab:vus} lists the
ten highest-priority candidates, all in \textit{MYH7}, with predicted
pathogenicity $>0.994$. At a threshold of $\hat{p}\geq 0.90$, a
substantial fraction of VUS are flagged as high-priority for
functional validation -- representing a targeted triage layer that
focuses expensive laboratory resources on highest-yield candidates.

\begin{table}[h]
\centering
\caption{Top-10 Prioritised VUS Candidates Ranked by Two-Tower
Predicted Pathogenicity. All listed variants carry current ClinVar
VUS status.}
\label{tab:vus}
\resizebox{\columnwidth}{!}{%
\begin{tabular}{lllc}
\toprule
\textbf{Gene} & \textbf{Variant} & \textbf{Reported Context} &
  $\hat{p}^{\,\mathrm{cal}}$ \\
\midrule
\textit{MYH7} & Y624D & Hypertrophic cardiomyopathy & 0.9960 \\
\textit{MYH7} & Y715C & Hypertrophic cardiomyopathy & 0.9953 \\
\textit{MYH7} & C672R & Adenocarcinoma context      & 0.9949 \\
\textit{MYH7} & F548C & Hypertrophic cardiomyopathy & 0.9948 \\
\textit{MYH7} & L329P & Hypertrophic cardiomyopathy & 0.9948 \\
\textit{MYH7} & D522V & Dilated cardiomyopathy      & 0.9945 \\
\textit{MYH7} & L396P & Hypertrophic cardiomyopathy & 0.9945 \\
\textit{MYH7} & Y624N & Hypertrophic cardiomyopathy & 0.9944 \\
\textit{MYH7} & F286C & Hypertrophic cardiomyopathy & 0.9944 \\
\textit{MYH7} & F456S & Cardiomyopathy              & 0.9943 \\
\bottomrule
\end{tabular}%
}
\end{table}


% ==============================================================
% SECTION VIII  --  DISCUSSION
% ==============================================================
\section{Discussion}
\label{sec:discussion}

\subsection{The Value of Disease-Specific Language Modelling}

Our results establish that a disease-specific, zero-leakage framework
substantially outperforms proteome-wide predictors across most
sarcomeric HCM genes. The two most informative comparisons are with
EVE and AlphaMissense: both represent state-of-the-art unsupervised
and semi-supervised approaches applied across the human proteome,
yet our Baseline RF exceeds them by 9.5 and 8.6 mean AUPRC points,
respectively, under identical LOGO evaluation conditions. The
performance advantage persists despite the fact that EVE and
AlphaMissense have access to far larger training corpora: the entire
metazoan and human proteome in the case of EVE, and the full human
proteome with AlphaFold structural context in the case of
AlphaMissense. This demonstrates that the disease-specific alignment
of training data, feature engineering, and evaluation protocol
provides a qualitatively different type of information than
proteome-scale generalisation.

The ESM-2 ablation result provides useful insights. The ESM-only
variant of the Two-Tower model achieves AUPRC 0.940 on
held-out \textit{TNNT2}, outperforming even the full fusion model
(0.931) and the Baseline RF (0.921). This implies that the 650M-
parameter evolutionary language model already encodes a superset of
the sarcomeric pathogenicity signal available in 51 handcrafted
biophysical descriptors. The handcrafted features do not provide
additive information in the simple concatenation fusion scheme, and
may introduce mild optimisation noise relative to the clean ESM-only
signal. This finding motivates more sophisticated fusion architectures -- 
cross-modal attention or gated modality weighting -- as productive
future directions.

\subsection{Addressing the Two Leakage Channels}

Prior computational evaluations of variant pathogenicity predictors
have rarely disentangled feature-level and split-level leakage
effects simultaneously. Our protocol enforces both in a single
experimental design: (i) removal of six leakage-prone feature groups
from the input matrix, and (ii) LOGO cross-validation that enforces
complete gene-level isolation. This combination has a measurable cost
in apparent performance relative to random-split evaluation, but
yields performance estimates with substantially higher deployment
relevance.

The performance gap between REVEL/MetaRNN and our models on
\textit{MYH7} and \textit{TNNI3} ($\Delta\mathrm{AUPRC}\approx 0.07$
in favour of external tools) is directly attributable to ClinVar
label saturation in these genes. \textit{MYH7} and \textit{TNNI3}
are the two most densely annotated sarcomeric genes in ClinVar,
providing a rich training signal to tools that incorporate population
frequency and prior classification evidence. Under LOGO conditions
that strip our models of such signal, the remaining molecular
discriminator must be purely sequence-based and structural. That our
models achieve AUPRC $>0.90$ on these genes without ACMG-derived
inputs confirms that the biophysical and evolutionary representations
carry genuine predictive content beyond database consensus.

\subsection{Translational Impact of ISM Heatmaps}

The ISM landscapes bridge the gap between statistical pathogenicity
prediction and mechanistic hypothesis generation. For a molecular
diagnostics laboratory or inherited cardiomyopathy clinic, the
clinical utility of an ISM heatmap is threefold. First, it provides
a \textit{residue-level triage}: positions that are ISM-hot across
all substitution classes represent structural bottlenecks where
functional validation assays should be prioritised. Second, it
enables \textit{panel-informed design}: saturation mutagenesis
experiments can be targeted to the specific substitution classes
predicted to be most destabilising, reducing assay cost. Third,
and most importantly for VUS re-classification, ISM-predicted
pathogenicity constitutes \emph{independent computational evidence}
under ACMG/AMP criterion PP3 (pathogenic supporting): because our
model is built without ACMG-derived inputs, its elevated prediction
represents mechanistic signal rather than restated clinical consensus.

\subsection{Understanding the Performance of the Random Forest vs. the Neural Network}

A notable finding from the internal LOGO validation is that the Baseline Random Forest outperformed the Two-Tower Hybrid Neural Network on raw discrimination (mean AUPRC 0.886 vs. 0.857) across seven of the nine panel genes. While neural architectures are frequently assumed to be superior for embedding-based tasks, this result aligns with well-documented phenomena in high-dimensional, low-sample-size regimes. Our dataset comprises 1,954 labeled variants represented in a 1,331-dimensional feature space. In this data regime, tree-based ensemble methods like Random Forests are exceptionally robust against overfitting, naturally performing feature selection through axis-aligned splits.

Furthermore, the strong performance of the RF model indicates that the frozen ESM-2 representations are already highly linearly separable with respect to pathogenicity. The Two-Tower model, attempting to learn a complex, cross-modal joint representation from heterogeneous data types (dense tabular constraints vs. high-dimensional semantic spaces), likely encounters mild optimization challenges or overfitting, despite the inclusion of dropout and batch normalization. However, as demonstrated by the Expected Calibration Error (ECE) metrics, the neural model achieves vastly superior probability calibration via Platt scaling, making its outputs significantly more reliable for downstream clinical thresholding than the uncalibrated probability smoothing of the tree ensemble. 

\subsection{Limitations}

We identify four principal limitations. First, the ESM-2 backbone
is frozen: fine-tuning on sarcomeric sequences could sharpen the
mutation perturbation signal but requires careful leakage control
if HCM-labelled sequences are included in the fine-tuning corpus.
Second, the framework does not incorporate explicit 3-D structural
coordinates; geometric deep learning over AlphaFold-derived contact
maps could capture long-range allosteric effects not fully
represented by sequence-level embeddings. Third, our training labels
are derived from ClinVar and may embed historical annotation errors
or inter-laboratory adjudication heterogeneity. Fourth, the LOGO
evaluation protocol, while more rigorous than random splits, still
evaluates within-species generalisation; extension to orthologous
non-human sarcomeric gene panels would further test biological
generalisability. Prospective validation against newly reclassified
variants accruing in ClinVar from 2025 onward is essential before
clinical deployment.


% ==============================================================
% SECTION IX  --  CONCLUSION
% ==============================================================
\section{Conclusion}
\label{sec:conclusion}

We have presented a comprehensive zero-leakage framework for
missense variant pathogenicity prediction in hypertrophic
cardiomyopathy that simultaneously addresses feature-level and
split-level data leakage. Two complementary classifiers -- a Baseline
Random Forest and a Two-Tower Hybrid Neural Network fusing 51
handcrafted biophysical descriptors with differential ESM-2
evolutionary embeddings -- are trained on 1,954 high-confidence labeled
sarcomeric variants and evaluated under strict Leave-One-Gene-Out
cross-validation. The Baseline RF achieves a mean AUPRC of
\textbf{0.886}, outperforming all five external benchmarks (EVE,
AlphaMissense, REVEL, MetaRNN, CardioBoost) in seven of nine target
genes. Platt-scaled Two-Tower probabilities yield a Brier score of
\textbf{0.145} on held-out \textit{TNNT2}, supporting clinical
interpretability. ISM heatmaps generated without any structural
supervision reproduce the biologically expected vulnerability
landscape of sarcomeric proteins, including the myosin converter
domain and troponin regulatory interface, demonstrating that the
learned representations encode genuine mechanistic knowledge rather
than statistical artifact.

A primary contribution of this work is the
demonstration that leakage-free, disease-specific representation
learning outperforms both general-purpose evolutionary tools and the
prior HCM-specific classifier in a homology-safe evaluation regime.
This constitutes a principled shift in the standard of evidence
required for computational variant pathogenicity claims: from
database consensus recapitulation to mechanism-grounded inference
validated by out-of-gene generalisation. The resulting probabilistic
risk rankings for 4,524 sarcomeric VUS provide an immediately
actionable triage resource for inherited cardiomyopathy clinics and
molecular diagnostics laboratories.


%  -- - Data Availability  --  --  --  --  --  --  --  --  --  --  --  --  -- -
\section*{Data and Code Availability}

All processed datasets, trained model weights, evaluation scripts,
and ISM generation code used in this study are available in the
associated project repository. The ESM-2 backbone is obtained from
the publicly available Hugging Face model hub at
\texttt{facebook/esm2\_t33\_650M\_UR50D}. CardioBoost precomputed
scores were downloaded from \url{https://cardiodb.org/cardioboost/}.
REVEL and MetaRNN scores were obtained from dbNSFP v4.5.


% ==============================================================
% REFERENCES
% ==============================================================
\begin{thebibliography}{25}

\bibitem{fowler2014deep}
D.~M. Fowler and S.~Fields,
``Deep mutational scanning: a new style of protein science,''
\textit{Nat. Methods}, vol.~11, no.~8, pp.~801--807, 2014.

\bibitem{ng2003sift}
P.~C. Ng and S.~Henikoff,
``SIFT: predicting amino acid changes that affect protein function,''
\textit{Nucleic Acids Res.}, vol.~31, no.~13, pp.~3812--3814, 2003.

\bibitem{adzhubei2010}
I.~A. Adzhubei \textit{et al.},
``A method and server for predicting damaging missense mutations,''
\textit{Nat. Methods}, vol.~7, no.~4, pp.~248--249, 2010.

\bibitem{richards2015}
S.~Richards \textit{et al.},
``Standards and guidelines for the interpretation of sequence variants:
a joint consensus recommendation of the ACMG and the AMP,''
\textit{Genet. Med.}, vol.~17, no.~5, pp.~405--424, 2015.

\bibitem{maron2013}
B.~J. Maron and M.~S. Maron,
``Hypertrophic cardiomyopathy,''
\textit{Lancet}, vol.~381, no.~9862, pp.~242--255, 2013.

\bibitem{semsarian2015}
C.~Semsarian \textit{et al.},
``New perspectives on the prevalence of hypertrophic cardiomyopathy,''
\textit{J.\ Am.\ Coll.\ Cardiol.}, vol.~65, no.~12,
pp.~1249--1254, 2015.

\bibitem{seidman2001}
C.~E. Seidman and J.~G. Seidman,
``The genetic basis for cardiomyopathy: from mutation identification
to mechanistic paradigms,''
\textit{Cell}, vol.~104, no.~4, pp.~557--567, 2001.

\bibitem{landrum2016clinvar}
M.~J. Landrum \textit{et al.},
``ClinVar: public archive of interpretations of clinically relevant
variants,''
\textit{Nucleic Acids Res.}, vol.~44, no.~D1,
pp.~D862--D868, 2016.

\bibitem{landrum2018clinvar}
M.~J. Landrum \textit{et al.},
``ClinVar: improving access to variant interpretations and supporting
evidence,''
\textit{Nucleic Acids Res.}, vol.~46, no.~D1,
pp.~D1062--D1067, 2018.

\bibitem{karczewski2020}
K.~J. Karczewski \textit{et al.},
``The mutational constraint spectrum quantified from variation in
141,456 humans,''
\textit{Nature}, vol.~581, no.~7809, pp.~434--443, 2020.

\bibitem{uniprot2023}
The UniProt Consortium,
``UniProt: the universal protein knowledgebase in 2023,''
\textit{Nucleic Acids Res.}, vol.~51, no.~D1,
pp.~D523--D531, 2023.

\bibitem{frazer2021eve}
J.~Frazer \textit{et al.},
``Disease variant prediction with deep generative models of
evolutionary data,''
\textit{Nature}, vol.~599, no.~7883, pp.~91--95, 2021.

\bibitem{cheng2023alphamissense}
J.~Cheng \textit{et al.},
``Accurate proteome-wide missense variant effect prediction with
AlphaMissense,''
\textit{Science}, vol.~381, no.~6664, p.~eadg7492, 2023.

\bibitem{ioannidis2016revel}
N.~M. Ioannidis \textit{et al.},
``REVEL: an ensemble method for predicting the pathogenicity of
rare missense variants,''
\textit{Am.\ J.\ Hum.\ Genet.}, vol.~99, no.~4,
pp.~877--885, 2016.

\bibitem{li2022metarnn}
C.~Li \textit{et al.},
``MetaRNN: differentiating rare pathogenic and rare benign missense
SNVs and InDels using deep learning,''
\textit{Genome Med.}, vol.~14, no.~1, p.~115, 2022.

\bibitem{zhang2021cardioboost}
X.~Zhang \textit{et al.},
``Disease-specific variant pathogenicity prediction significantly
improves variant interpretation in inherited cardiac conditions,''
\textit{Genet. Med.}, vol.~23, no.~1, pp.~69--79, 2021.

\bibitem{lin2023esm2}
Z.~Lin \textit{et al.},
``Evolutionary-scale prediction of atomic-level protein structure
with a language model,''
\textit{Science}, vol.~379, no.~6637, pp.~1123--1130, 2023.

\bibitem{grimm2015circularity}
D.~G. Grimm \textit{et al.},
``The evaluation of tools used to predict the impact of missense
variants is hindered by two types of circularity,''
\textit{Hum.\ Mutat.}, vol.~36, no.~5, pp.~513--523, 2015.

\bibitem{kapoor2023leakage}
S.~Kapoor and A.~Narayanan,
``Leakage and the reproducibility crisis in machine-learning-based
science,''
\textit{Patterns}, vol.~4, no.~9, p.~100804, 2023.

\bibitem{breiman2001rf}
L.~Breiman,
``Random forests,''
\textit{Mach. Learn.}, vol.~45, no.~1, pp.~5--32, 2001.

\bibitem{loshchilov2019adamw}
I.~Loshchilov and F.~Hutter,
``Decoupled weight decay regularization,''
in \textit{Proc. ICLR}, 2019.

\bibitem{spudich2019}
J.~A. Spudich,
``Three perspectives on the molecular basis of hypercontractility
caused by hypertrophic cardiomyopathy mutations,''
\textit{Pfl\"{u}gers Arch.}, vol.~471, no.~5, pp.~701--717, 2019.

\bibitem{grantham1974}
R.~Grantham,
``Amino acid difference formula to help explain protein evolution,''
\textit{Science}, vol.~185, no.~4154, pp.~862--864, 1974.

\bibitem{nag2021}
S.~Nag \textit{et al.},
``Contractility parameters of human $\beta$-cardiac myosin
with the hypertrophic cardiomyopathy mutation R403Q show loss of
motor function,''
\textit{Sci.\ Adv.}, vol.~1, no.~9, p.~e1500511, 2015.

\bibitem{platt1999}
J.~C. Platt,
``Probabilistic outputs for support vector machines and comparisons
to regularized likelihood methods,''
in \textit{Advances in Large Margin Classifiers},
A.~J. Smola \textit{et al.}, eds., MIT Press, 1999, pp.~61--74.

\end{thebibliography}

\end{document}
